\chapter{Swirl a vortice libero}
\label{ch:swirl}

Tutto quanto precede ha assunto che il gas si muova in un piano
contenente l'asse: in avanti e verso l'esterno, mai attorno. Un
motore a detonazione rotante (RDE) --- un combustore in cui un'onda
di detonazione corre di continuo attorno a una camera anulare --- non
è così accomodante: il suo scarico esce \emph{ruotando}. Questo
capitolo aggiunge quella rotazione al solutore nell'unica forma che
non costa nuove ipotesi, mostra per dimostrazione e per misura che
l'intera macchina certificata le sopravvive, e dichiara con chiarezza
che cosa l'estensione modella e che cosa no.

\section{La domanda}

Poiché l'onda di detonazione spazza la camera, trascina con sé il
gas: lo scarico porta una \emph{velocità azimutale} $w$ --- la
componente di velocità che gira \emph{attorno} all'asse,
perpendicolare al piano passante per l'asse in cui il solutore lavora
(quel piano si chiama piano \emph{meridiano}, e $q_m$ indica la
velocità in esso). Il corpus della teoria garantisce (la sua
affermazione N6-2) che per un particolare profilo di swirl --- il
\emph{vortice libero} --- l'intera struttura della marcia certificata
sopravvive intatta, e che qualunque cosa oltre esso è una classe di
flusso che questo metodo non può rappresentare affatto, e che
appartiene alla linea dell'aggiunto CFD.

Un vortice libero è la rotazione dello scarico di una vasca lontano
dal suo nucleo: il prodotto del raggio per la velocità azimutale,
$\Gamma = y\,w$, è \emph{uno e lo stesso numero ovunque nel flusso}.
Quel prodotto si chiama \emph{circolazione}; misura quanto il gas
gira attorno all'asse e fisicamente è il momento angolare per unità
di massa. Lo swirl è dunque più veloce vicino all'asse e si spegne
verso l'esterno, ed è fissato da un solo scalare $\Gamma$.

La domanda di questo capitolo: il motore sa portare lo swirl, con la
garanzia \emph{misurata} anziché assunta --- e con lo swirl spento
($\Gamma \to 0$) che recupera esattamente la marcia certificata senza
swirl?

\section{Come funziona, per esteso}

\subsection{Perché il vortice libero è l'unica estensione gratuita}

Tre fatti esatti rendono speciale il vortice libero.

Primo, l'equazione della quantità di moto attorno all'asse dice che
$y\,w$ è trasportato invariato lungo ciascuna linea di corrente. Se è
una sola costante \emph{ovunque}, quell'equazione è soddisfatta
identicamente e $w = \Gamma/y$ è noto dalla sola posizione ---
nessuna nuova incognita entra nella marcia.

Secondo, il flusso resta \emph{irrotazionale}: le particelle di
fluido vengono trasportate e deformate, ma non ruotano attorno al
proprio centro. È questa la proprietà su cui posa l'intero solutore
--- è ciò che rende il flusso descrivibile con le due famiglie
caratteristiche del Capitolo~\ref{ch:moc}; un flusso le cui
particelle ruotano si dice \emph{rotazionale}, e per esso il metodo
delle caratteristiche non è più applicabile. È esattamente ciò che
sta oltre il vortice libero. Il teorema di Crocco, che lega la
vorticità ai gradienti di entalpia totale ed entropia, lascia
possibile solo una componente azimutale quando entrambe sono
uniformi, e la condizione $V \times \omega = 0$ costringe poi anche
quella ad annullarsi. Il flusso meridiano resta dunque irrotazionale,
e la struttura certificata a due caratteristiche non è qui
un'approssimazione ma è esatta.

Terzo, svolgendo le equazioni della gasdinamica con $w = \Gamma/y$ si
vede che il problema meridiano cambia in esattamente due punti e in
nessun altro:
\begin{itemize}[itemsep=1pt]
\item \textbf{lo stato}: l'energia cinetica della rotazione entra nel
  bilancio energetico, $h = h_0 - q_m^2/2 - \Gamma^2/(2y^2)$, quindi
  lo stato termodinamico dipende ora dal raggio $y$ oltre che dalla
  velocità meridiana $q_m$;
\item \textbf{la sorgente}: un gas che ruota dev'essere trattenuto da
  un gradiente di pressione diretto verso l'interno, e quel
  contributo centrifugo si somma al termine che il flusso
  assialsimmetrico già porta --- il \emph{termine sorgente
  assialsimmetrico}, il pezzo aggiuntivo nelle relazioni di
  compatibilità che tiene conto dell'area di flusso che si allarga al
  crescere di $y$, assente nel flusso piano. Diventa
  $c^2 v/y \;\to\; (c^2 + \Gamma^2/y^2)\,v/y$.
\end{itemize}
Le direzioni caratteristiche mantengono la loro forma, ora
all'angolo di Mach \emph{meridiano}, e i coefficienti di
compatibilità in $(u, v)$ restano intatti.

\subsection{L'implementazione sono due modifiche, e soltanto due}

Il cambiamento di stato non costa nulla di nuovo. Con la velocità
aumentata $q_{\rm eff} = \sqrt{q_m^2 + \Gamma^2/y^2}$ la chiusura
energetica si legge $h = h_0 - q_{\rm eff}^2/2$, esattamente la forma
senza swirl, quindi la macchina termodinamica certificata si applica
alla lettera a $q_{\rm eff}$ e solo il numero di Mach viene
ricostruito dalla velocità meridiana. Il bordo di getto libero --- la
frontiera dove lo scarico incontra l'atmosfera, sulla quale la
pressione è assegnata --- mantiene la sua parametrizzazione
certificata: $p = p_a$ fissa la velocità \emph{totale}, lo stesso
scalare che senza swirl, quindi la velocità meridiana del bordo segue
una legge esplicita, $q_m(y) = \sqrt{q_{p_a}^2 - \Gamma^2/y^2}$.

I tre tipi di cella sono rispecchiati con lo stato e la sorgente di
swirl e poi eseguiti attraverso il driver di marcia certificato
\emph{invariato} tramite una giunzione additiva, i cui valori
predefiniti riproducono esattamente il percorso certificato ---
entrambe le suite della marcia del plug sono state rieseguite dopo la
giunzione e continuano a passare.

È emersa una lezione numerica, ora corretta nelle celle di swirl. Il
piccolo sistema non lineare della cella mescola righe di dimensioni
molto diverse: righe geometriche di ordine uno contro righe di
compatibilità di ordine $u^2$. L'iterazione di Newton smorzata ---
che accetta un passo solo se una norma del residuo diminuisce ---
può allora arenarsi a un passo dalla radice, perché il passo che
sistema la posizione baratta un aumento minuscolo nelle righe di
compatibilità contro una correzione geometrica di ordine $10^{-7}$, e
lo smorzamento lo rigetta per sempre. Dividere le righe di
compatibilità per una velocità di riferimento costante --- un puro
riscalamento delle equazioni che non muove alcuna soluzione ---
ripristina la convergenza certificata (cella peggiore $0.009$ della
tolleranza).

\subsection{Che cosa questo è e che cosa non è: il motivo di fase a
spirale}
\label{sec:spiral}

Il vortice libero modella la velocità azimutale \emph{media nel
tempo} dello scarico. Non modella lo \emph{svolgersi a spirale del
motivo di fase}: in un RDE reale le superfici a fase costante sono
eliche --- in un dato istante, stazioni assiali diverse (e azimut
diversi) si trovano in punti diversi del ciclo, perché l'onda ruota
mentre il flusso transita. Lo strato di ciclo del
Capitolo~\ref{ch:cycle} è deliberatamente \emph{quasi-stazionario}:
ogni fase riempie l'intero ugello come flusso stazionario, e il ciclo
è una media sulle fasi senza orologio. Quella chiusura mette da parte
due effetti fisici distinti. Primo, la \emph{non stazionarietà entro
un raggio meridiano}: la fase vista alla stazione $x$ è in ritardo
sulla fase all'iniettore del tempo di transito, e se il periodo del
ciclo è confrontabile con il tempo che il gas passa nell'ugello, il
flusso non si equilibra mai ad alcuna fase --- il modello
quasi-stazionario è il limite di frequenza ridotta nulla, e per
numeri rappresentativi di un RDE (rotazione a kHz contro tempi di
residenza sub-millisecondo) quella frequenza è di ordine uno, quindi
questo è un limite reale del modello, dichiarato. Secondo,
l'\emph{accoppiamento azimutale}: la disposizione elicoidale inclina
i fronti d'onda fuori dal piano meridiano --- gasdinamica
genuinamente tridimensionale che nessun modello assialsimmetrico per
fase può vedere. Entrambi gli effetti portano con sé una mitigazione
parziale e una via d'uscita pulita. La mitigazione: se i raggi
meridiani fossero disaccoppiati e quasi-stazionari, la \emph{media
temporale} in ogni stazione eguaglierebbe la media sulle fasi
indipendentemente da come le fasi sono disposte nello spazio --- la
spirale rimescola il \emph{quando}, non il \emph{che cosa};
l'obiettivo mediato sul ciclo è insensibile alla disposizione a
quell'ordine. La via d'uscita: entrambe le deviazioni vivono
esattamente dove il corpus già indica --- la linea CFD non
stazionaria tridimensionale, le cui simulazioni contengono
nativamente la spirale e per le quali questo motore è la banca degli
oracoli.

\subsection{I mondi}

Due mondi di prova, entrambi sulle tabelle di gas reale, entrambi con
uno swirl forte ($w/u = 0.35$ al raggio interno).

Il \emph{condotto a equilibrio radiale} è una risposta nota esatta.
In un condotto anulare rettilineo la soluzione con swirl ha velocità
assiale \emph{uniforme} e velocità radiale nulla: derivando la
chiusura energetica lungo $y$ e sottraendo il bilancio radiale fra
gradiente di pressione e forza centrifuga,
$\mathrm{d}p/\mathrm{d}y = \rho\Gamma^2/y^3$, si ottiene
$\mathrm{d}u/\mathrm{d}y = 0$ esattamente --- mentre l'intera
struttura radiale di pressione, temperatura, velocità del suono e
numero di Mach è portata dal termine di swirl
(Figura~\ref{fig:swirl}a). La marcia deve \emph{conservare} quello
stato, e ogni deriva è puro errore.

Il \emph{mondo in espansione} fa scendere dolcemente la parete sotto
la stessa corrente, così il flusso è genuinamente bidimensionale
($v \neq 0$) e la sorgente centrifuga agisce davvero. Il solo
condotto non può vedere quella sorgente, perché lì $v = 0$ la
annulla.

\section{Che cosa è stato costruito}

\carrier{validation/a1\_swirl\_march.py}: lo stato con swirl
(velocità aumentata) e i suoi coefficienti, i tre processi di cella
rispecchiati con il riscalamento delle righe, gli integrali di flusso
e di spinta di parete consapevoli dello swirl, i due mondi e la suite
qui sotto. Il driver certificato ha acquisito la giunzione additiva
(\code{cells=}, \code{q\_edge=}); le sue suite di oracolo e
assialsimmetrica sono state rieseguite 6/6 e 5/5 dopo la modifica.

\section{Risultati di riferimento: 9/9 PASS}

W-1, la regressione a swirl spento: a circolazione nulla le celle di
swirl riproducono la marcia delle celle certificate a meno di
$2.1\times10^{-13}$, contro una banda di $10^{-9}$ relativi --- la
tolleranza certificata di Newton per cella, il che significa che il
sistema non lineare di ogni cella è stato risolto a un residuo dentro
una banda derivata. L'unica differenza aritmetica è la radice
quadrata della velocità aumentata.

W-2, il condotto esatto: tutte le celle certificate da Newton
(peggiore $0.009$ della tolleranza), e l'equilibrio è conservato allo
zero di macchina --- deriva della velocità assiale
$1.6\times10^{-15}$, $|v|/u\ 2.3\times10^{-15}$, pressione di parete
$4.9\times10^{-15}$, posizione del bordo esatta --- dentro bande di
Richardson, gli errori ammessi ottenuti rieseguendo con la spaziatura
delle stazioni dimezzata e lasciando che lo schema misuri il proprio
errore di troncamento (Capitolo~\ref{ch:method}).

W-3, il mondo in espansione: la massa chiude a $7.8$ su
$4.3\times10^{4}$ (banda $16$), e il teorema della quantità di moto
in direzione assiale chiude in banda --- invariato dallo swirl,
perché la forza centrifuga è radiale e la quantità di moto attorno
all'asse ne è disaccoppiata; verificato, non assunto.

W-4, i due controlli di vitalità: spegnere lo swirl sposta la
velocità di parete $8.6\times$ oltre la banda propria della coppia, e
lasciar cadere \emph{solo} la sorgente centrifuga mantenendo $\Gamma$
nello stato la sposta $12\times$ oltre --- quindi il nuovo termine è
genuinamente collegato e genuinamente rilevante
(Figura~\ref{fig:swirl}b). R-1, il rigettatore: celle alimentate con
una circolazione corrotta ($\Gamma^2 \times 1.10$) ma marciate su
dati di equilibrio veri escono dalle bande di conservazione di dodici
ordini di grandezza --- colte, come richiesto.

\begin{figure}[htbp]
\centering
\includegraphics[width=\textwidth]{fig_swirl.pdf}
\caption{Swirl a vortice libero, misurato. (a)~Il condotto a
equilibrio radiale, una risposta nota esatta: raggio $y$ [m] in
verticale, pressione statica $p$ [bar] e numero di Mach meridiano $M$
[--] in orizzontale (due scale). La pressione sale e il numero di
Mach cala verso il raggio esterno benché la velocità assiale sia
uniforme e la velocità radiale sia nulla --- l'intera struttura
radiale è portata dal termine di swirl nel bilancio energetico. La
marcia conserva questo stato a $10^{-15}$. (b)~Il mondo in
espansione: coordinata assiale $x$ [m] contro pressione di parete
$p_w$ [bar], per tre esecuzioni --- swirl completo, swirl con la
sorgente centrifuga deliberatamente eliminata, e nessuno swirl. Le
tre curve si separano molto oltre le bande d'errore ammesse
($8.6\times$ e $12\times$), il che dimostra che sia lo swirl sia il
suo termine sorgente sono collegati e rilevanti. (c)~Che cosa il
vortice libero modella davvero: raggio $y$ [m] contro velocità
azimutale $w = \Gamma/y$ [m/s] --- la rotazione media nel tempo dello
scarico anulare, più veloce al raggio interno e in decadimento verso
l'esterno. È una \emph{media}, non il motivo di fase a spirale del
motore reale, che è fisica tridimensionale non stazionaria e
appartiene alla linea CFD (Sezione~\ref{sec:spiral}).}
\label{fig:swirl}
\end{figure}
